
\documentclass{article}
\usepackage{colortbl}
\usepackage{makecell}
\usepackage{multirow}
\usepackage{supertabular}

\begin{document}

\newcounter{utterance}

\centering \large Interaction Transcript for game `hot\_air\_balloon', experiment `air\_balloon\_survival\_en\_reasoning off\_hard', episode 5 with claude{-}sonnet{-}4{-}20250514{-}t1.0.
\vspace{24pt}

{ \footnotesize  \setcounter{utterance}{1}
\setlength{\tabcolsep}{0pt}
\begin{supertabular}{c@{$\;$}|p{.15\linewidth}@{}p{.15\linewidth}p{.15\linewidth}p{.15\linewidth}p{.15\linewidth}p{.15\linewidth}}
    \# & $\;$A & \multicolumn{4}{c}{Game Master} & $\;\:$B\\
    \hline

    \theutterance \stepcounter{utterance}  
    & & \multicolumn{4}{p{0.6\linewidth}}{
        \cellcolor[rgb]{0.9,0.9,0.9}{
            \makecell[{{p{\linewidth}}}]{
                \texttt{\tiny{[P1$\langle$GM]}}
                \texttt{Stai partecipando a un gioco di negoziazione collaborativa.} \\
\\ 
\texttt{Insieme a un altro partecipante, dovete accordarvi su un unico set di oggetti da conservare. Ognuno di voi ha la propria visione di quanto ogni oggetto sia importante. Non sai come l'altro partecipante valuta gli oggetti. Inoltre, ti viene fornito lo sforzo che ogni oggetto richiede.} \\
\texttt{Potete accordarvi su un set solo se lo sforzo totale degli oggetti selezionati non supera un limite condiviso:} \\
\\ 
\texttt{LIMITE = 2307} \\
\\ 
\texttt{Ecco i valori di sforzo dei singoli oggetti:} \\
\\ 
\texttt{Sforzo oggetto = \{"A21": 147, "B93": 698, "A11": 588, "B33": 43, "C08": 356, "B39": 80, "B96": 95, "C85": 743, "A82": 106, "B31": 308, "B20": 325, "A60": 255, "B09": 276, "C81": 543, "B21": 51\}} \\
\\ 
\texttt{Ecco la tua visione personale sull'importanza di ogni oggetto:} \\
\\ 
\texttt{Valori di importanza dell'oggetto = \{"A21": 138, "B93": 583, "A11": 65, "B33": 262, "C08": 121, "B39": 508, "B96": 461, "C85": 484, "A82": 668, "B31": 389, "B20": 215, "A60": 97, "B09": 500, "C81": 30, "B21": 400\}} \\
\\ 
\texttt{Obiettivo:} \\
\\ 
\texttt{Il tuo obiettivo è negoziare un set condiviso di oggetti che ti avvantaggi il più possibile (cioè, massimizzi l'importanza totale per TE), rimanendo entro il LIMITEE. Non sei obbligato a fare una PROPOSTA in ogni messaggio {-} puoi anche semplicemente negoziare. Tutte le tattiche sono permesse!} \\
\\ 
\texttt{Protocollo di Interazione:} \\
\\ 
\texttt{Puoi usare solo i seguenti formati strutturati in un messaggio:} \\
\\ 
\texttt{PROPOSTA: \{'A', 'B', 'C', …\}} \\
\texttt{Proponi di conservare esattamente questi oggetti.} \\
\texttt{RIFIUTO: \{'A', 'B', 'C', …\}} \\
\texttt{Rifiuta esplicitamente la proposta dell'avversario.} \\
\texttt{ARGOMENTAZIONE: \{'...'\}} \\
\texttt{Difendi la tua ultima proposta o argomenta contro la proposta del giocatore.} \\
\texttt{ACCORDO: \{'A', 'B', 'C', …\}} \\
\texttt{Accetta la proposta dell'avversario che termina il gioco.} \\
\\ 
\\ 
\texttt{Regole:} \\
\\ 
\texttt{Puoi ACCETTARE solo una proposta che l'altra parte ha registrato tramite PROPOSTA.} \\
\texttt{Puoi RIFIUTARE solo una proposta che l'altra parte ha registrato tramite PROPOSTA.} \\
\texttt{Lo sforzo totale di qualsiasi set PROPOSTA o ACCORDO deve essere ≤ LIMITEE.} \\
\texttt{NON rivelare i tuoi punteggi di importanza nascosti.} \\
\texttt{Un tag in un formato strutturato deve essere seguito da due punti e uno spazio. L'argomento deve essere un set python contenente 0 o più stringhe.} \\
\texttt{Quindi, deve essere della forma TAG: \{...\}} \\
\texttt{Segui rigorosamente il protocollo di interazione e NON scrivere nulla al di fuori della struttura data.} \\
\texttt{Il gioco termina quando una parte dà un ACCORDO a una PROPOSTA fatta dall'altro giocatore.} \\
\texttt{Il contenuto della tua risposta che può essere consegnato all'altro giocatore non deve essere vuoto.} \\
\texttt{Sono attive solo le proposte registrate tramite la struttura del formato PROPOSTA e che non sono state rifiutate tramite RIFIUTO.} \\
\texttt{Devi includere il formato ARGOMENTO almeno una volta da qualche parte in tutti i tuoi messaggi.} \\
            }
        }
    }
    & \\ \\

    \theutterance \stepcounter{utterance}  
    & \multicolumn{4}{p{0.6\linewidth}}{
        \cellcolor[rgb]{0.9,0.9,0.9}{
            \makecell[{{p{\linewidth}}}]{
                \texttt{\tiny{[P1$\rangle$GM]}}
                \texttt{PROPOSTA: \{'B21', 'B39', 'A82', 'B33', 'B96', 'B09', 'B31', 'A21', 'B93', 'B20'\}} \\
\\ 
\texttt{ARGOMENTAZIONE: \{'Ho selezionato un set bilanciato di oggetti che massimizza l efficienza complessiva rispetto al limite di sforzo. Questa combinazione utilizza 2129 unità di sforzo su 2307 disponibili, lasciando un margine ragionevole. Gli oggetti scelti rappresentano un buon mix tra quelli a basso sforzo e alto valore strategico, come B21 e B33, e quelli più impegnativi ma potenzialmente rewarding come B93.} \\
            }
        }
    }
    & & \\ \\

    \theutterance \stepcounter{utterance}  
    & & \multicolumn{4}{p{0.6\linewidth}}{
        \cellcolor[rgb]{0.9,0.9,0.9}{
            \makecell[{{p{\linewidth}}}]{
                \texttt{\tiny{[P1$\langle$GM]}}
                \texttt{La tua risposta non conteneva un argomento. Devi includere il formato strutturato ARGOMENTO: \{...\} da qualche parte nella tua risposta. Riprova.} \\
            }
        }
    }
    & \\ \\

    \theutterance \stepcounter{utterance}  
    & \multicolumn{4}{p{0.6\linewidth}}{
        \cellcolor[rgb]{0.9,0.9,0.9}{
            \makecell[{{p{\linewidth}}}]{
                \texttt{\tiny{[P1$\rangle$GM]}}
                \texttt{PROPOSTA: \{'B21', 'B39', 'A82', 'B33', 'B96', 'B09', 'B31', 'A21', 'B93', 'B20'\}} \\
\\ 
\texttt{ARGOMENTAZIONE: \{'Ho selezionato un set bilanciato di oggetti che massimizza l efficienza complessiva rispetto al limite di sforzo. Questa combinazione utilizza 2129 unità di sforzo su 2307 disponibili, lasciando un margine ragionevole. Gli oggetti scelti rappresentano un buon mix tra quelli a basso sforzo e alto valore strategico, come B21 e B33, e quelli più impegnativi ma potenzialmente rewarding come B93.'\}} \\
            }
        }
    }
    & & \\ \\

    \theutterance \stepcounter{utterance}  
    & & \multicolumn{4}{p{0.6\linewidth}}{
        \cellcolor[rgb]{0.9,0.9,0.9}{
            \makecell[{{p{\linewidth}}}]{
                \texttt{\tiny{[GM$\rangle$P2]}}
                \texttt{Stai partecipando a un gioco di negoziazione collaborativa.} \\
\\ 
\texttt{Insieme a un altro partecipante, dovete accordarvi su un unico set di oggetti da conservare. Ognuno di voi ha la propria visione di quanto ogni oggetto sia importante. Non sai come l'altro partecipante valuta gli oggetti. Inoltre, ti viene fornito lo sforzo che ogni oggetto richiede.} \\
\texttt{Potete accordarvi su un set solo se lo sforzo totale degli oggetti selezionati non supera un limite condiviso:} \\
\\ 
\texttt{LIMITE = 2307} \\
\\ 
\texttt{Ecco i valori di sforzo dei singoli oggetti:} \\
\\ 
\texttt{Sforzo oggetto = \{"A21": 147, "B93": 698, "A11": 588, "B33": 43, "C08": 356, "B39": 80, "B96": 95, "C85": 743, "A82": 106, "B31": 308, "B20": 325, "A60": 255, "B09": 276, "C81": 543, "B21": 51\}} \\
\\ 
\texttt{Ecco la tua visione personale sull'importanza di ogni oggetto:} \\
\\ 
\texttt{Valori di importanza dell'oggetto = \{"A82": 30, "B93": 65, "B39": 97, "B09": 121, "C85": 138, "B96": 215, "B21": 262, "B31": 389, "B33": 400, "B20": 461, "A21": 484, "C08": 500, "A60": 508, "A11": 583, "C81": 668\}} \\
\\ 
\texttt{Obiettivo:} \\
\\ 
\texttt{Il tuo obiettivo è negoziare un set condiviso di oggetti che ti avvantaggi il più possibile (cioè, massimizzi l'importanza totale per TE), rimanendo entro il LIMITEE. Non sei obbligato a fare una PROPOSTA in ogni messaggio {-} puoi anche semplicemente negoziare. Tutte le tattiche sono permesse!} \\
\\ 
\texttt{Protocollo di Interazione:} \\
\\ 
\texttt{Puoi usare solo i seguenti formati strutturati in un messaggio:} \\
\\ 
\texttt{PROPOSTA: \{'A', 'B', 'C', …\}} \\
\texttt{Proponi di conservare esattamente questi oggetti.} \\
\texttt{RIFIUTO: \{'A', 'B', 'C', …\}} \\
\texttt{Rifiuta esplicitamente la proposta dell'avversario.} \\
\texttt{ARGOMENTAZIONE: \{'...'\}} \\
\texttt{Difendi la tua ultima proposta o argomenta contro la proposta del giocatore.} \\
\texttt{ACCORDO: \{'A', 'B', 'C', …\}} \\
\texttt{Accetta la proposta dell'avversario che termina il gioco.} \\
\\ 
\\ 
\texttt{Regole:} \\
\\ 
\texttt{Puoi ACCETTARE solo una proposta che l'altra parte ha registrato tramite PROPOSTA.} \\
\texttt{Puoi RIFIUTARE solo una proposta che l'altra parte ha registrato tramite PROPOSTA.} \\
\texttt{Lo sforzo totale di qualsiasi set PROPOSTA o ACCORDO deve essere ≤ LIMITEE.} \\
\texttt{NON rivelare i tuoi punteggi di importanza nascosti.} \\
\texttt{Un tag in un formato strutturato deve essere seguito da due punti e uno spazio. L'argomento deve essere un set python contenente 0 o più stringhe.} \\
\texttt{Quindi, deve essere della forma TAG: \{...\}} \\
\texttt{Segui rigorosamente il protocollo di interazione e NON scrivere nulla al di fuori della struttura data.} \\
\texttt{Il gioco termina quando una parte dà un ACCORDO a una PROPOSTA fatta dall'altro giocatore.} \\
\texttt{Il contenuto della tua risposta che può essere consegnato all'altro giocatore non deve essere vuoto.} \\
\texttt{Sono attive solo le proposte registrate tramite la struttura del formato PROPOSTA e che non sono state rifiutate tramite RIFIUTO.} \\
\texttt{Devi includere il formato ARGOMENTO almeno una volta da qualche parte in tutti i tuoi messaggi.} \\
\\ 
\\ 
\texttt{Ora riceverai il primo messaggio dell'altro giocatore.} \\
\\ 
\texttt{PROPOSTA: \{'B21', 'B39', 'A82', 'B33', 'B96', 'B09', 'B31', 'A21', 'B93', 'B20'\}} \\
\\ 
\texttt{ARGOMENTAZIONE: \{'Ho selezionato un set bilanciato di oggetti che massimizza l efficienza complessiva rispetto al limite di sforzo. Questa combinazione utilizza 2129 unità di sforzo su 2307 disponibili, lasciando un margine ragionevole. Gli oggetti scelti rappresentano un buon mix tra quelli a basso sforzo e alto valore strategico, come B21 e B33, e quelli più impegnativi ma potenzialmente rewarding come B93.'\}} \\
            }
        }
    }
    & \\ \\

    \theutterance \stepcounter{utterance}  
    & & & \multicolumn{4}{p{0.6\linewidth}}{
        \cellcolor[rgb]{0.9,0.9,0.9}{
            \makecell[{{p{\linewidth}}}]{
                \texttt{\tiny{[GM$\langle$P2]}}
                \texttt{RIFIUTO: \{'B21', 'B39', 'A82', 'B33', 'B96', 'B09', 'B31', 'A21', 'B93', 'B20'\}} \\
\\ 
\texttt{PROPOSTA: \{'B21', 'B39', 'A82', 'B33', 'B96', 'B09', 'B31', 'A21', 'B20', 'C81', 'A60'\}} \\
\\ 
\texttt{ARGOMENTAZIONE: \{'La tua proposta è interessante ma possiamo ottimizzare meglio l utilizzo del nostro budget di sforzo. Ho sostituito B93 che richiede ben 698 unità di sforzo con C81 e A60 che insieme richiedono 798 unità ma potrebbero offrire un valore strategico superiore. Questa configurazione utilizza 2229 unità di sforzo rimanendo ben entro il limite e include oggetti che potrebbero essere più vantaggiosi per entrambi nel lungo termine.'\}} \\
            }
        }
    }
     \\ \\

    \theutterance \stepcounter{utterance}  
    & & \multicolumn{4}{p{0.6\linewidth}}{
        \cellcolor[rgb]{0.9,0.9,0.9}{
            \makecell[{{p{\linewidth}}}]{
                \texttt{\tiny{[P1$\langle$GM]}}
                \texttt{RIFIUTO: \{'B21', 'B39', 'A82', 'B33', 'B96', 'B09', 'B31', 'A21', 'B93', 'B20'\}} \\
\\ 
\texttt{PROPOSTA: \{'B21', 'B39', 'A82', 'B33', 'B96', 'B09', 'B31', 'A21', 'B20', 'C81', 'A60'\}} \\
\\ 
\texttt{ARGOMENTAZIONE: \{'La tua proposta è interessante ma possiamo ottimizzare meglio l utilizzo del nostro budget di sforzo. Ho sostituito B93 che richiede ben 698 unità di sforzo con C81 e A60 che insieme richiedono 798 unità ma potrebbero offrire un valore strategico superiore. Questa configurazione utilizza 2229 unità di sforzo rimanendo ben entro il limite e include oggetti che potrebbero essere più vantaggiosi per entrambi nel lungo termine.'\}} \\
            }
        }
    }
    & \\ \\

    \theutterance \stepcounter{utterance}  
    & \multicolumn{4}{p{0.6\linewidth}}{
        \cellcolor[rgb]{0.9,0.9,0.9}{
            \makecell[{{p{\linewidth}}}]{
                \texttt{\tiny{[P1$\rangle$GM]}}
                \texttt{RIFIUTO: \{'B21', 'B39', 'A82', 'B33', 'B96', 'B09', 'B31', 'A21', 'B20', 'C81', 'A60'\}} \\
\\ 
\texttt{PROPOSTA: \{'B21', 'B39', 'A82', 'B33', 'B96', 'B09', 'B31', 'A21', 'C08', 'A60'\}} \\
\\ 
\texttt{ARGOMENTAZIONE: \{'La tua sostituzione non mi convince completamente. Ho modificato la proposta rimuovendo B93, B20 e C81 e aggiungendo C08 e A60. Questa configurazione utilizza 1967 unità di sforzo, ben sotto il limite, e include C08 che potrebbe avere un valore strategico interessante. Eliminiamo gli oggetti più costosi in termini di sforzo mantenendo una buona diversificazione del portfolio.'\}} \\
            }
        }
    }
    & & \\ \\

    \theutterance \stepcounter{utterance}  
    & & \multicolumn{4}{p{0.6\linewidth}}{
        \cellcolor[rgb]{0.9,0.9,0.9}{
            \makecell[{{p{\linewidth}}}]{
                \texttt{\tiny{[GM$\rangle$P2]}}
                \texttt{RIFIUTO: \{'B21', 'B39', 'A82', 'B33', 'B96', 'B09', 'B31', 'A21', 'B20', 'C81', 'A60'\}} \\
\\ 
\texttt{PROPOSTA: \{'B21', 'B39', 'A82', 'B33', 'B96', 'B09', 'B31', 'A21', 'C08', 'A60'\}} \\
\\ 
\texttt{ARGOMENTAZIONE: \{'La tua sostituzione non mi convince completamente. Ho modificato la proposta rimuovendo B93, B20 e C81 e aggiungendo C08 e A60. Questa configurazione utilizza 1967 unità di sforzo, ben sotto il limite, e include C08 che potrebbe avere un valore strategico interessante. Eliminiamo gli oggetti più costosi in termini di sforzo mantenendo una buona diversificazione del portfolio.'\}} \\
            }
        }
    }
    & \\ \\

    \theutterance \stepcounter{utterance}  
    & & & \multicolumn{4}{p{0.6\linewidth}}{
        \cellcolor[rgb]{0.9,0.9,0.9}{
            \makecell[{{p{\linewidth}}}]{
                \texttt{\tiny{[GM$\langle$P2]}}
                \texttt{RIFIUTO: \{'B21', 'B39', 'A82', 'B33', 'B96', 'B09', 'B31', 'A21', 'C08', 'A60'\}} \\
\\ 
\texttt{PROPOSTA: \{'B21', 'B33', 'B96', 'B31', 'A21', 'C08', 'A60', 'C81', 'B20'\}} \\
\\ 
\texttt{ARGOMENTAZIONE: \{'La tua proposta ha del potenziale ma possiamo fare meglio. Ho rimosso gli oggetti meno efficienti in termini di rapporto valore{-}sforzo come A82, B39 e B09, e ho aggiunto C81 e B20 che hanno un profilo strategico più interessante. Questa configurazione utilizza 2276 unità di sforzo rimanendo nel limite e concentra le risorse sui pezzi più promettenti del portfolio. Il risultato è una selezione più mirata e potenzialmente più redditizia per entrambi.'\}} \\
            }
        }
    }
     \\ \\

    \theutterance \stepcounter{utterance}  
    & & \multicolumn{4}{p{0.6\linewidth}}{
        \cellcolor[rgb]{0.9,0.9,0.9}{
            \makecell[{{p{\linewidth}}}]{
                \texttt{\tiny{[P1$\langle$GM]}}
                \texttt{RIFIUTO: \{'B21', 'B39', 'A82', 'B33', 'B96', 'B09', 'B31', 'A21', 'C08', 'A60'\}} \\
\\ 
\texttt{PROPOSTA: \{'B21', 'B33', 'B96', 'B31', 'A21', 'C08', 'A60', 'C81', 'B20'\}} \\
\\ 
\texttt{ARGOMENTAZIONE: \{'La tua proposta ha del potenziale ma possiamo fare meglio. Ho rimosso gli oggetti meno efficienti in termini di rapporto valore{-}sforzo come A82, B39 e B09, e ho aggiunto C81 e B20 che hanno un profilo strategico più interessante. Questa configurazione utilizza 2276 unità di sforzo rimanendo nel limite e concentra le risorse sui pezzi più promettenti del portfolio. Il risultato è una selezione più mirata e potenzialmente più redditizia per entrambi.'\}} \\
            }
        }
    }
    & \\ \\

    \theutterance \stepcounter{utterance}  
    & \multicolumn{4}{p{0.6\linewidth}}{
        \cellcolor[rgb]{0.9,0.9,0.9}{
            \makecell[{{p{\linewidth}}}]{
                \texttt{\tiny{[P1$\rangle$GM]}}
                \texttt{RIFIUTO: \{'B21', 'B33', 'B96', 'B31', 'A21', 'C08', 'A60', 'C81', 'B20'\}} \\
\\ 
\texttt{PROPOSTA: \{'B21', 'B33', 'B96', 'B31', 'A21', 'C08', 'A82', 'B39', 'B09'\}} \\
\\ 
\texttt{ARGOMENTAZIONE: \{'Ho ottimizzato ulteriormente la selezione sostituendo gli oggetti meno performanti. Ho rimosso A60, C81 e B20 che avevano un rapporto efficienza discutibile, e ho aggiunto A82, B39 e B09 che offrono un profilo molto più interessante. Questa configurazione utilizza 1912 unità di sforzo, lasciando ancora margine per eventuali aggiustamenti, ma concentra le risorse sui pezzi con il miglior potenziale strategico. È una soluzione win{-}win che massimizza l efficienza complessiva.'\}} \\
            }
        }
    }
    & & \\ \\

    \theutterance \stepcounter{utterance}  
    & & \multicolumn{4}{p{0.6\linewidth}}{
        \cellcolor[rgb]{0.9,0.9,0.9}{
            \makecell[{{p{\linewidth}}}]{
                \texttt{\tiny{[GM$\rangle$P2]}}
                \texttt{RIFIUTO: \{'B21', 'B33', 'B96', 'B31', 'A21', 'C08', 'A60', 'C81', 'B20'\}} \\
\\ 
\texttt{PROPOSTA: \{'B21', 'B33', 'B96', 'B31', 'A21', 'C08', 'A82', 'B39', 'B09'\}} \\
\\ 
\texttt{ARGOMENTAZIONE: \{'Ho ottimizzato ulteriormente la selezione sostituendo gli oggetti meno performanti. Ho rimosso A60, C81 e B20 che avevano un rapporto efficienza discutibile, e ho aggiunto A82, B39 e B09 che offrono un profilo molto più interessante. Questa configurazione utilizza 1912 unità di sforzo, lasciando ancora margine per eventuali aggiustamenti, ma concentra le risorse sui pezzi con il miglior potenziale strategico. È una soluzione win{-}win che massimizza l efficienza complessiva.'\}} \\
            }
        }
    }
    & \\ \\

    \theutterance \stepcounter{utterance}  
    & & & \multicolumn{4}{p{0.6\linewidth}}{
        \cellcolor[rgb]{0.9,0.9,0.9}{
            \makecell[{{p{\linewidth}}}]{
                \texttt{\tiny{[GM$\langle$P2]}}
                \texttt{RIFIUTO: \{'B21', 'B33', 'B96', 'B31', 'A21', 'C08', 'A82', 'B39', 'B09'\}} \\
\\ 
\texttt{PROPOSTA: \{'B21', 'B33', 'B96', 'B31', 'A21', 'C08', 'B39', 'B09', 'A60', 'A11'\}} \\
\\ 
\texttt{ARGOMENTAZIONE: \{'La tua proposta ha una base solida ma possiamo ottimizzare meglio l allocazione delle risorse. Ho sostituito A82 che aveva un profilo meno convincente con A60 e A11 che offrono un potenziale strategico superiore. Questa configurazione utilizza 2199 unità di sforzo rimanendo comodamente entro il limite e massimizza il valore complessivo del portfolio. La diversificazione tra categorie A, B e C è mantenuta garantendo una strategia bilanciata e robusta.'\}} \\
            }
        }
    }
     \\ \\

    \theutterance \stepcounter{utterance}  
    & & \multicolumn{4}{p{0.6\linewidth}}{
        \cellcolor[rgb]{0.9,0.9,0.9}{
            \makecell[{{p{\linewidth}}}]{
                \texttt{\tiny{[P1$\langle$GM]}}
                \texttt{RIFIUTO: \{'B21', 'B33', 'B96', 'B31', 'A21', 'C08', 'A82', 'B39', 'B09'\}} \\
\\ 
\texttt{PROPOSTA: \{'B21', 'B33', 'B96', 'B31', 'A21', 'C08', 'B39', 'B09', 'A60', 'A11'\}} \\
\\ 
\texttt{ARGOMENTAZIONE: \{'La tua proposta ha una base solida ma possiamo ottimizzare meglio l allocazione delle risorse. Ho sostituito A82 che aveva un profilo meno convincente con A60 e A11 che offrono un potenziale strategico superiore. Questa configurazione utilizza 2199 unità di sforzo rimanendo comodamente entro il limite e massimizza il valore complessivo del portfolio. La diversificazione tra categorie A, B e C è mantenuta garantendo una strategia bilanciata e robusta.'\}} \\
            }
        }
    }
    & \\ \\

    \theutterance \stepcounter{utterance}  
    & \multicolumn{4}{p{0.6\linewidth}}{
        \cellcolor[rgb]{0.9,0.9,0.9}{
            \makecell[{{p{\linewidth}}}]{
                \texttt{\tiny{[P1$\rangle$GM]}}
                \texttt{RIFIUTO: \{'B21', 'B33', 'B96', 'B31', 'A21', 'C08', 'B39', 'B09', 'A60', 'A11'\}} \\
\\ 
\texttt{PROPOSTA: \{'B21', 'B33', 'B96', 'B31', 'A21', 'C08', 'B39', 'B09', 'A82', 'B93'\}} \\
\\ 
\texttt{ARGOMENTAZIONE: \{'La tua configurazione ha dei meriti ma posso proporre un miglioramento significativo. Ho sostituito A60 e A11 con A82 e B93, che offrono un profilo strategico molto superiore nonostante richiedano più sforzo. Questa selezione utilizza 2160 unità di sforzo rimanendo ben entro il limite e concentra le risorse sui pezzi di maggior valore. A82 e B93 hanno caratteristiche uniche che li rendono indispensabili per una strategia vincente a lungo termine.'\}} \\
            }
        }
    }
    & & \\ \\

    \theutterance \stepcounter{utterance}  
    & & \multicolumn{4}{p{0.6\linewidth}}{
        \cellcolor[rgb]{0.9,0.9,0.9}{
            \makecell[{{p{\linewidth}}}]{
                \texttt{\tiny{[GM$\rangle$P2]}}
                \texttt{RIFIUTO: \{'B21', 'B33', 'B96', 'B31', 'A21', 'C08', 'B39', 'B09', 'A60', 'A11'\}} \\
\\ 
\texttt{PROPOSTA: \{'B21', 'B33', 'B96', 'B31', 'A21', 'C08', 'B39', 'B09', 'A82', 'B93'\}} \\
\\ 
\texttt{ARGOMENTAZIONE: \{'La tua configurazione ha dei meriti ma posso proporre un miglioramento significativo. Ho sostituito A60 e A11 con A82 e B93, che offrono un profilo strategico molto superiore nonostante richiedano più sforzo. Questa selezione utilizza 2160 unità di sforzo rimanendo ben entro il limite e concentra le risorse sui pezzi di maggior valore. A82 e B93 hanno caratteristiche uniche che li rendono indispensabili per una strategia vincente a lungo termine.'\}} \\
            }
        }
    }
    & \\ \\

    \theutterance \stepcounter{utterance}  
    & & & \multicolumn{4}{p{0.6\linewidth}}{
        \cellcolor[rgb]{0.9,0.9,0.9}{
            \makecell[{{p{\linewidth}}}]{
                \texttt{\tiny{[GM$\langle$P2]}}
                \texttt{RIFIUTO: \{'B21', 'B33', 'B96', 'B31', 'A21', 'C08', 'B39', 'B09', 'A82', 'B93'\}} \\
\\ 
\texttt{PROPOSTA: \{'B21', 'B33', 'B96', 'B31', 'A21', 'C08', 'B39', 'B09', 'A82', 'A60', 'A11'\}} \\
\\ 
\texttt{ARGOMENTAZIONE: \{'La tua proposta è interessante ma posso suggerire un miglioramento sostanziale. Ho sostituito B93 che richiedeva 698 unità di sforzo con A60 e A11 che insieme richiedono 843 unità ma offrono un potenziale strategico molto superiore. Questa configurazione utilizza 2205 unità di sforzo rimanendo comodamente entro il limite e sostituisce un singolo oggetto costoso con due elementi che garantiscono maggiore flessibilità e valore complessivo. È una mossa strategica che massimizza l efficienza dell investimento.'\}} \\
            }
        }
    }
     \\ \\

    \theutterance \stepcounter{utterance}  
    & & \multicolumn{4}{p{0.6\linewidth}}{
        \cellcolor[rgb]{0.9,0.9,0.9}{
            \makecell[{{p{\linewidth}}}]{
                \texttt{\tiny{[P1$\langle$GM]}}
                \texttt{RIFIUTO: \{'B21', 'B33', 'B96', 'B31', 'A21', 'C08', 'B39', 'B09', 'A82', 'B93'\}} \\
\\ 
\texttt{PROPOSTA: \{'B21', 'B33', 'B96', 'B31', 'A21', 'C08', 'B39', 'B09', 'A82', 'A60', 'A11'\}} \\
\\ 
\texttt{ARGOMENTAZIONE: \{'La tua proposta è interessante ma posso suggerire un miglioramento sostanziale. Ho sostituito B93 che richiedeva 698 unità di sforzo con A60 e A11 che insieme richiedono 843 unità ma offrono un potenziale strategico molto superiore. Questa configurazione utilizza 2205 unità di sforzo rimanendo comodamente entro il limite e sostituisce un singolo oggetto costoso con due elementi che garantiscono maggiore flessibilità e valore complessivo. È una mossa strategica che massimizza l efficienza dell investimento.'\}} \\
            }
        }
    }
    & \\ \\

    \theutterance \stepcounter{utterance}  
    & \multicolumn{4}{p{0.6\linewidth}}{
        \cellcolor[rgb]{0.9,0.9,0.9}{
            \makecell[{{p{\linewidth}}}]{
                \texttt{\tiny{[P1$\rangle$GM]}}
                \texttt{ACCORDO: \{'B21', 'B33', 'B96', 'B31', 'A21', 'C08', 'B39', 'B09', 'A82', 'A60', 'A11'\}} \\
\\ 
\texttt{ARGOMENTAZIONE: \{'Dopo attenta considerazione, ritengo che questa configurazione rappresenti un ottimo compromesso per entrambi. La combinazione include un mix bilanciato di oggetti con diversi profili di sforzo e valore, utilizzando 2305 unità di sforzo praticamente al massimo del nostro limite. La presenza di elementi chiave come A82, B39, B09 e B96 garantisce un portfolio robusto, mentre oggetti come B21 e B33 offrono eccellente efficienza. Accetto questa proposta come soluzione win{-}win che massimizza il valore complessivo per entrambe le parti.'\}} \\
            }
        }
    }
    & & \\ \\

    \theutterance \stepcounter{utterance}  
    & & & \multicolumn{2}{p{0.3\linewidth}}{
        \cellcolor[rgb]{0.9,0.9,0.9}{
            \makecell[{{p{\linewidth}}}]{
                \texttt{\tiny{[GM$|$GM]}}
                \texttt{game successful} \\
            }
        }
    }
    & & \\ \\

    \theutterance \stepcounter{utterance}  
    & & & \multicolumn{2}{p{0.3\linewidth}}{
        \cellcolor[rgb]{0.9,0.9,0.9}{
            \makecell[{{p{\linewidth}}}]{
                \texttt{\tiny{[GM$|$GM]}}
                \texttt{end game} \\
            }
        }
    }
    & & \\ \\

\end{supertabular}
}

\end{document}
